\chapter{Literature Review and Related Work}
\label{chap:relatedworks}

\section{Competitor Analysis}
\label{section:competitor-analysis}

\textbf{Manual CCTV Monitoring Systems} are the most widely deployed approach in many cities, including Thailand. Traffic officers review footage to identify violations, making this method labor-intensive, error-prone, and impossible to scale city-wide.

\textbf{Standalone ANPR Systems} from vendors such as Genetec and Hikvision are mature commercial products designed for vehicle identification, including toll collection, parking, and stolen-vehicle detection. They do not perform helmet detection and are not designed for motorcycle-specific workflows.

\textbf{Academic research} has demonstrated that AI can detect helmet violations in traffic footage, but these systems were built solely as experiments to test the concept. Most only focus on detecting the helmet and do not include license plate recognition, violation logging, or any real-world deployment on actual roads.

\textbf{HelmTrack} addresses the gap left by all three categories by combining real-time helmet detection and license plate recognition into a single automated pipeline, eliminating the need for manual review while producing structured, evidence-backed violation records.

\section{Literature Review}
\label{section:literature-review}

<TIP: Review academic papers or technical reports
that are relevant to your project.
For each source, briefly summarize its key findings and
then explain how it informs your own design or implementation decisions.
The goal is not just to describe what others have done,
but to build a justification for your own technical choices.
For example, if you chose a particular algorithm or architecture,
cite the work that supports that decision. />
